\chapter{分数与幂}\label{ch:g9:fractions}

分数与幂是计算中的常用工具，可用于分配
数量、比较比例、书写很大或很小的
数。本章整理分数与幂的运算规则 ---
每个中间步骤都写出来 ---
并引入科学记数法。

\section{分数的运算}

\begin{definition}[分数]\label{def:g9:fractions:fraction}
形如 $\dfrac ab$（$b \neq 0$）的数称为\emph{分数}\index{分数}，它是
即 $a$ 除以 $b$ 的商。若两分数中一个可由
另一个通过将分子与分母同乘（或同除）
以同一非零数得到，则两分数相等：
\[
\frac ab = \frac{a \times k}{b \times k} \qquad (k \neq 0).
\]
\end{definition}

\begin{example}\label{ex:g9:fractions:simplify}
逐步化简 $\dfrac{42}{56}$：
\[
\frac{42}{56} = \frac{42 \div 2}{56 \div 2} = \frac{21}{28}
= \frac{21 \div 7}{28 \div 7} = \frac34 .
\]
（在\cref{ch:g9:arith}中，最大公因数
可以一步完成这一化简。）
\end{example}

\begin{method}[分数的加减法]\label{met:g9:fractions:add}
\begin{enumerate}
  \item 将两个分数通分，取\emph{公分母}（两个分母的
        公倍数）；
  \item 将分子相加或相减，分母保持不变；
  \item 若可能，将结果化简。
\end{enumerate}
\end{method}

\begin{example}\label{ex:g9:fractions:add}
计算 $\dfrac56 + \dfrac34$。$6$ 与 $4$ 的一个公分母是
$12$：
\[
\frac56 + \frac34
= \frac{5 \times 2}{6 \times 2} + \frac{3 \times 3}{4 \times 3}
= \frac{10}{12} + \frac{9}{12}
= \frac{19}{12}.
\]
计算 $2 - \dfrac37$。将 $2$ 写成分母为 $7$ 的分数：
\[
2 - \frac37 = \frac{14}{7} - \frac37 = \frac{11}{7}.
\]
\end{example}

\begin{omfigure}
\begin{tikzpicture}[scale=0.55]
  \foreach \i in {0,...,11}
    \draw[omExo] ({mod(\i,12)},0) rectangle ({mod(\i,12)+1},1);
  \foreach \i in {0,...,9}
    \fill[omDef!35] (\i,0) rectangle (\i+1,1);
  \draw[thick] (0,0) rectangle (12,1);
  \node[below, font=\small] at (6,-0.1)
    {$\frac56 = \frac{10}{12}$：12 格中的 10 格};
  \begin{scope}[yshift=-2.4cm]
  \foreach \i in {0,...,11}
    \draw[omExo] (\i,0) rectangle (\i+1,1);
  \foreach \i in {0,...,8}
    \fill[omThm!35] (\i,0) rectangle (\i+1,1);
  \draw[thick] (0,0) rectangle (12,1);
  \node[below, font=\small] at (6,-0.1)
    {$\frac34 = \frac{9}{12}$：12 格中的 9 格};
  \end{scope}
\end{tikzpicture}

{\small 为何要通分：将两个整体各切成 $12$ 个
等大的格子，这样两个分数便于比较、便于相加 ---
合计为 $10 + 9 = 19$ 个十二分之一。}
\end{omfigure}

\begin{method}[分数的乘除法]\label{met:g9:fractions:mult}
\begin{enumerate}
  \item 相乘时，分别将分子与分子、分母与分母
        相乘：$\dfrac ab \times \dfrac cd = \dfrac{ac}{bd}$
        （只要可能，应\emph{先约分}再相乘）；
  \item 相除时，乘以除数的倒数：
        $\dfrac ab \div \dfrac cd = \dfrac ab \times \dfrac dc$.
\end{enumerate}
\end{method}

\begin{example}\label{ex:g9:fractions:mult}
\[
\frac{7}{15} \times \frac{25}{14}
= \frac{7 \times 25}{15 \times 14}
= \frac{7 \times 25}{15 \times 2 \times 7}
= \frac{25}{30} = \frac56,
\]
其中先约去公因数 $7$，再约去 $5$。再看一个
除法的例子：
\[
\frac{3}{4} \div \frac{9}{8}
= \frac34 \times \frac89
= \frac{3 \times 8}{4 \times 9}
= \frac{24}{36} = \frac23 .
\]
\end{example}

\section{幂}

\begin{definition}[幂]\label{def:g9:fractions:power}
对实数 $a$ 与正整数 $n$，
\emph{幂}\index{幂} $a^n$ 是 $n$ 个等于
$a$ 的因数之积：
\[
a^n = \underbrace{a \times a \times \dots \times a}_{n \text{ 个因数}} .
\]
按约定，$a^0 = 1$（$a \neq 0$），负指数
表示倒数：
\[
a^{-n} = \frac{1}{a^n}.
\]
\end{definition}

\begin{example}
$2^4 = 2 \times 2 \times 2 \times 2 = 16$; \quad
$(-3)^2 = 9$ 而 $-3^2 = -(3^2) = -9$（指数的运算先于
负号）；\quad $10^{-3} = \frac{1}{1000} = 0.001$； \quad
$5^1 = 5$ 和 $5^0 = 1$。
\end{example}

\begin{theorem}[指数运算法则]\label{thm:g9:fractions:rules}
对任意非零的 $a$、$b$ 及所有整数 $m$、$n$：
\[
a^m \times a^n = a^{m+n},
\qquad
\frac{a^m}{a^n} = a^{m-n},
\qquad
\left(a^m\right)^n = a^{mn},
\qquad
(ab)^n = a^n b^n .
\]
\end{theorem}

\begin{proof}
正指数的情形，数一数因数的个数。对于第一条法则：
$a^m \times a^n$ 是 $m$ 个因数 $a$ 后接 $n$ 个因数 $a$，共
$m + n$ 个因数。对于第三条：
$\left(a^m\right)^n$ 表示 $n$ 组各含 $m$ 个因数的积，得到
$mn$ 个因数。其余两条法则类似；零指数与负指数的情形
由 $a^{-n} = \frac{1}{a^n}$ 检验（并使
约定 $a^0 = 1$ 成为唯一相容的选择，因为
$a^n \times a^0$ 必须等于 $a^{n+0} = a^n$）。
\end{proof}

\begin{example}\label{ex:g9:fractions:rules}
逐步化简：
\[
\frac{3^7 \times 3^{-2}}{3^4}
= \frac{3^{7 + (-2)}}{3^4}
= \frac{3^5}{3^4}
= 3^{5-4} = 3^1 = 3 .
\]
再看含字母的情形：$(2a^3)^2 \times a = 4 a^6 \times a = 4a^7$。
\end{example}

\section{十的幂与科学记数法}

\begin{proposition}[十的幂]\label{prop:g9:fractions:ten}
对正整数 $n$，$10^n$ 写作“$1$ 后接 $n$
个零”，且 $10^{-n} = 0.0\dots01$ 中的 $1$ 位于小数点后第 $n$
位。小数乘以 $10^n$（或 $10^{-n}$）时，其小数点
向右（或向左）移动 $n$ 位。
\end{proposition}

\begin{definition}[科学记数法]\label{def:g9:fractions:scientific}
正小数的\emph{科学记数法}\index{科学记数法}
是将其唯一地写成下列形式：
\[
a \times 10^n,
\qquad\text{其中 } 1 \leq a < 10 \text{ 且 } n \text{ 为整数。}
\]
\end{definition}

\begin{example}\label{ex:g9:fractions:scientific}
地球与太阳的距离约为 $149\,600\,000$ km $= 1.496 \times
10^8$ km。一个水分子的大小约为 $0.000\,000\,000\,28$ m
$= 2.8 \times 10^{-10}$ m。计算这类数时，
分别结合小数部分与十的幂：
\[
(3 \times 10^5) \times (2.5 \times 10^{-2})
= (3 \times 2.5) \times 10^{5 + (-2)}
= 7.5 \times 10^{3}.
\]
\end{example}

\begin{method}[将数写成科学记数法]\label{met:g9:fractions:scientific}
\begin{enumerate}
  \item 将小数点移到第一个非零数字之后，这样
        便得到因数 $a$，满足 $1 \leq a < 10$；
  \item 数出小数点移动的位数，这个位数即为 $n$，
        原数 $\geq 10$ 时取正，原数
        $< 1$ 时取负；
  \item 检验：$a \times 10^n$ 须还原原数。
\end{enumerate}
\end{method}

\section{练习}

\begin{exercise}[$\star$]\label{exo:g9:fractions:1}
计算，并将结果化为最简分数：
\[
\frac13 + \frac15, \qquad
\frac72 - \frac53, \qquad
\frac{5}{12} + \frac38, \qquad
3 - \frac45 .
\]
\end{exercise}

\begin{exercise}[$\star$]\label{exo:g9:fractions:2}
计算并化简：
\[
\frac59 \times \frac{27}{35}, \qquad
\frac{8}{15} \div \frac{4}{5}, \qquad
\frac23 \times \frac34 \times \frac45 .
\]
\end{exercise}

\begin{exercise}[$\star$]\label{exo:g9:fractions:3}
计算：
\[
2^5, \qquad (-2)^4, \qquad -2^4, \qquad 4^{-2}, \qquad
\left(\tfrac23\right)^3, \qquad 7^0 .
\]
\end{exercise}

\begin{exercise}[$\star$]\label{exo:g9:fractions:4}
写成单个幂的形式：
\[
5^3 \times 5^6, \qquad
\frac{7^9}{7^5}, \qquad
\left(2^3\right)^4, \qquad
\frac{3^2 \times 3^7}{3^5}, \qquad
2^5 \times 4 .
\]
\end{exercise}

\begin{exercise}[$\star$]\label{exo:g9:fractions:5}
用科学记数法表示：$45\,000$；$0.0072$；$302.5$；
$0.000\,000\,9$；一千二百万。
\end{exercise}

\begin{exercise}[$\star\star$]\label{exo:g9:fractions:6}
计算，并用科学记数法表示结果：
\[
(2 \times 10^7) \times (4.5 \times 10^{-3}),
\qquad
\frac{9 \times 10^5}{3 \times 10^{-2}},
\qquad
(5 \times 10^3)^2 .
\]
\end{exercise}

\begin{exercise}[$\star\star$]\label{exo:g9:fractions:7}
水箱已注入容积的 $\frac35$。加入 $30$ 升水后，水箱注入容积的 $\frac34$。
这个水箱的容积是多少？（先将加入的水量
表示为水箱容积的分数。）
\end{exercise}

\begin{exercise}[$\star\star$]\label{exo:g9:fractions:8}
爱丽丝吃掉一块蛋糕的 $\frac14$，然后本吃掉剩下部分的 $\frac13$，
接着卡拉吃掉剩下部分的 $\frac12$。最后，整块蛋糕中
还剩下的分数是多少？（逐步算出每一步之后的剩余部分。）
\end{exercise}

\begin{exercise}[$\star\star$]\label{exo:g9:fractions:9}
光每秒传播 $3 \times 10^8$ 米。太阳与我们相距约
$1.5 \times 10^{11}$ 米。阳光需要多长时间才能到达我们？
请先以秒为单位给出答案（不要求科学记数法），再
以分钟和秒为单位给出。
\end{exercise}

\begin{exercise}[$\star\star\star$]\label{exo:g9:fractions:10}
哪一个更大，$2^{30}$ 还是 $3^{20}$？（提示：把两者都写成指数为 $10$ 的幂，
利用 $\left(a^m\right)^n = a^{mn}$，并比较 $2^3$ 与
$3^2$ 的大小。）
\end{exercise}

\section{问题：循环小数的秘密代码}

\begin{problem}[{周末问题 --- 每个分数的小数写法都会终止或循环，
每个循环小数都是一个分数，
且 $0.999\ldots$ 恰好等于 $1$}]
\label{pb:g9:fractions:1}
用 $3$ 除以 $11$，数字便陷入一段反复的吟唱：
$0.272727\ldots$ 无限延续。用 $1$ 除以 $4$，数字则戛然而止
停在 $0.25$。本题将证明，分数只有这两种
可能的小数行为，然后倒转这台机器：
给定任一循环小数，还原它背后隐藏的分数。
沿途它还将一劳永逸地解决学校数学中
最具争议的一个等式。

\medskip
\noindent\textbf{第一部分 --- 从分数到小数。}
\begin{enumerate}
  \item 用长除法求出
        $\frac14$、$\frac13$、$\frac56$ 与 $\frac{3}{11}$ 的小数写法。
        哪些终止，哪些循环，循环节是什么？
  \item 说明\emph{任一}分数
        $\frac ab$ 的小数写法必然终止或最终循环。（在
        用 $b$ 做长除法时，相继的余数可以取哪些值？
        某个余数出现的那一刻，会发生什么？
        尤其是它第二次出现的时候？）
  \item 将 $1$ 除以 $7$ 的除法进行到足够远，以找出
        循环节。它有多长？这一长度与
        第 2 小题所给出的界限相比如何？
  \item 有些分数之所以终止，是因为其分母
        可以扩大为 10 的幂：$\frac14 =
        \frac{25}{100}$, $\frac{3}{8} = \frac{375}{1000}$.
        说明这一判别准则，以及为什么没有任何整数乘子
        能把 $3$、$6$ 或 $11$ 变为 $10$、$100$，
        $1000, \dots$（乘积的个位数字会是什么？）。
  \item 不做除法，将 $\frac{9}{40}$、$\frac{5}{12}$ 与
        $\frac{33}{50}$ 分为“终止”与“循环”两类，再写出
        其中两个终止的分数的小数写法。
\end{enumerate}

\medskip
\noindent\textbf{第二部分 --- 从循环小数回到
分数。}
\begin{enumerate}[resume]
  \item 移位技巧。设 $x = 0.7777\ldots$，计算 $10x$，
        再计算 $10x - x$，由此求出 $x$ 的分数形式。用
        除法进行检验。
  \item 设 $x = 0.727272\ldots$。哪一个 10 的幂
        恰好使 $x$ 移动一个循环节？利用它把 $x$ 写成
        最简分数。
  \item 混合情形：$x = 0.58333\ldots$（只有 $3$
        循环）。计算 $1000x - 100x$，由此求出 $x$ 的
        最简分数。
  \item 最著名的例子：对
        $x = 0.9999\ldots$ 应用移位技巧，你得到的是什么分数——
        也就是哪一个\emph{数}——？请将结果与你的直觉相调和，
        记住\cref{pb:g6:decimals:1}中的“幽灵般的邻居”，
        以及\cref{pb:g6:fractions:1}中逐步缩小的余数：
        $0.999\ldots$ 是否“刚好小于”$1$，
        还是 $1$ 的另一种写法？
  \item 将 $2.454545\ldots$ 与 $0.123123123\ldots$ 化为
        最简分数。
\end{enumerate}

\medskip
\noindent\textbf{第三部分 --- 字典已补全。}
\begin{enumerate}[resume]
  \item 第 6--10 题提示了一部“字典”：
        从小数点起循环的 $k$ 位数字组成的循环节，等于这段数字
        为分子、以 $k$ 个 9 为分母的分数（$0.727272\ldots = \frac{72}{99}$）。用
        这部字典预测 $0.037037\ldots$ 对应的分数，
        并将其化简——再惊叹一下：哪个出人意料地简单的分数，
        其循环节恰为“$037$”？用除法检验你的结果。
  \item 将字典与移位结合，
        将 $0.0272727\ldots$ 化为最简分数。
  \item 数 $0.101001000100001\ldots$（先一个 $1$，再一个
        $0$，然后一个 $1$，接着两个 $0$，再一个 $1$，然后三个
        $0$，依此类推）永不停止。它是循环小数吗？
        由第一部分可知，它能否是
        分数？（你刚刚遇到了第一个可证明的
        \emph{无理数}；最著名的一个
        在\cref{ch:g9:sqrt}中出现。）
  \item 估计 $2^{30}$ 大约有多少位数字？请用
        $2^{10} = 1\,024 \approx 1.024 \times 10^3$ 与
        指数法则（\cref{thm:g9:fractions:rules}），将
        $2^{30}$ 写成（近似的）科学记数法，并
        作出结论。（比较\cref{exo:g9:fractions:10}。）
  \item 收尾：将 $0.20262026\ldots$ 化为分数，并
        陈述整个问题的双向结论：哪些小数写法是分数，
        而哪些分数具有怎样的小数写法，
        即各分数对应哪些小数写法？
\end{enumerate}
\end{problem}
